\section{模型的评价、部署闭环与推广}
\label{sec:evaluation}

\subsection{关键稳健性证据}

本文不再额外拼接一组统一的参数扰动试验，而是从各问已经完成的对照中选取三项最能影响结论的证据，汇总于表~\ref{tab:robustness-evidence}。Q1 的结果说明日内成本对可转移能量规模更敏感；Q3 的年末边界对照排除了通过年末透支改变策略排序的可能；Q4 则给出场景规模在成本精度与在线求解时间之间的折中。

\begin{table}[htbp]
\centering
\caption{关键稳健性证据汇总（均为已有试验结果）}
\label{tab:robustness-evidence}
\begingroup
\setlength{\tabcolsep}{5pt}
\renewcommand{\arraystretch}{1.28}
\begin{tabularx}{0.98\linewidth}{@{}C{1.0cm}C{3.1cm}YY@{}}
\toprule
问题 & 核验设置 & 已有结果 & 支持的结论 \\
\midrule
Q1 & 容量与功率上限分别变动 $\pm10\%$
& 容量变动后成本为 $34547.91/33092.98$ 元；功率变动后为 $33827.76/33780.88$ 元
& 同幅度变化下，容量对成本的影响大于充放电功率上限 \\
Q3 & 年末 SOC 由 $1200$ 调整为 $6000\,\mathrm{kWh}$
& 主策略成本增加 $2262.88$ 元，五种策略的相对排序保持不变
& 在已测试的两种年末边界下，策略排序保持不变 \\
Q4 & 在同一闭环协议下比较 $K=8$ 与 $K=12$
& $K=8$ 成本高 $39308.25$ 元（$0.24\%$），平均耗时由 $6.7$ s 降至 $5.6$ s，约缩短 $16.3\%$
& 在既定 $1\%$ 成本容差内，$K=8$ 以较小成本差换取更快求解 \\
\bottomrule
\end{tabularx}
\endgroup
\end{table}


这三项检验分别覆盖设备尺度、跨期边界与算法规模，结论均限定在题给数据、候选集合和当前结算解释之内，不外推为所有微网场景下的普适规律。

\subsection{面向部署的运行闭环}

若将本文策略实现为调度原型，运行链路可按以下三个环节闭合：

\begin{enumerate}[label=\textbf{D\arabic*：},leftmargin=4.2em,labelsep=0.4em,itemsep=0.3em]
\item \textbf{日初计划。}每日 0:00 接入截至当时的负荷、光伏、价格、预报和 BMS-SOC，生成并锁定普通购电额度或基准承诺。
\item \textbf{日内控制与安全门。}每 10 分钟以最新已观测状态重建剩余时域 LP，只执行下一步；6:00、12:00、18:00 仅可调整未执行承诺，且预测收益须覆盖改约成本。EMS 校验能量平衡、额度与弃光，BMS 校验 SOC、功率与互斥，任一硬约束不满足即拒绝下发。
\item \textbf{降级与留痕。}求解超时、数据缺失或通信中断时冻结已锁定承诺，储能转入 BMS 限幅安全保持，缺口由紧急购电补足；保存预测版本、信息截止点、承诺、SOC、求解状态和成本分项，确保费用与信息边界均可逐时复核。
\end{enumerate}

这一闭环把预测、优化、执行和结算分开记录：预测误差不会被隐藏为“求解失败”，控制动作也不能用事后信息改写，从而保留模型在真实运行环境中的可解释性与责任边界。

\subsection{工程边界与后续扩展}

题目未给出公共耦合点（PCC）的允许交换容量，因此本文主模型没有设置外网购电功率上限，当前结果应理解为\textbf{经济调度原型}，不能直接视为已经满足并网容量校核的现场方案。实际部署前应由配电设计或并网协议给出 $G_{\max}$，在每个时段加入
\begin{equation}
0\le g_t\le G_{\max}\Delta t,
\label{eq:pcc-limit-extension}
\end{equation}
并对 Q1--Q4 的可行性、紧急购电和成本重新校验；必要时还需加入变压器容量、线路潮流与备用约束。

本文也未将电池循环退化计入主模型。若能够获得单位吞吐退化成本，可在目标函数中加入
\begin{equation}
\sum_t c_{\mathrm{deg}}\bigl(c_t+r_t\bigr),
\label{eq:degradation-extension}
\end{equation}
其中 $c_{\mathrm{deg}}$ 为单位充放电吞吐量的折算退化成本。该项目前只是后续扩展，不能把现有成本结果解释为已经权衡电池寿命。类似地，需求响应、余电上网和多能源耦合也须在取得真实设备与市场参数后另行建模。
